% source: Robin Hartshorne, "Algebraic Geometry", GTM 52 (1977)
% pdf_page: 0242
% doc_page: 225 (Chapter III, Section 5, "The Cohomology of Projective Space")
% retrieved: 2026-07-11
% transcribed_by: claude-sonnet-5 (vision, rendered at 200dpi via PyMuPDF)

\DeclareMathOperator{\Proj}{Proj}

\section*{5 The Cohomology of Projective Space}

In this section we make explicit calculations of the cohomology of the
sheaves $\mathcal{O}(n)$ on a projective space, by using \v{C}ech cohomology
for a suitable open affine covering. These explicit calculations form the
basis for various general results about cohomology of coherent sheaves on
projective varieties.

Let $A$ be a noetherian ring, let $S = A[x_0, \ldots, x_r]$, and let
$X = \Proj S$ be the projective space $\mathbf{P}^r_A$ over $A$. Let
$\mathcal{O}_X(1)$ be the twisting sheaf of Serre (II, \S 5). For any sheaf
of $\mathcal{O}_X$-modules $\mathcal{F}$, we denote by $\Gamma_*(\mathcal{F})$
the graded $S$-module $\bigoplus_{n \in \mathbf{Z}} \Gamma(X, \mathcal{F}(n))$
(see II, \S 5).

\begin{theorem}[5.1]
Let $A$ be a noetherian ring, and let $X = \mathbf{P}^r_A$, with $r \ge 1$.
Then:
\begin{enumerate}
\item[(a)] the natural map $S \to \Gamma_*(\mathcal{O}_X) =
  \bigoplus_{n \in \mathbf{Z}} H^0(X, \mathcal{O}_X(n))$ is an isomorphism of
  graded $S$-modules;
\item[(b)] $H^i(X, \mathcal{O}_X(n)) = 0$ for $0 < i < r$ and all
  $n \in \mathbf{Z}$;
\item[(c)] $H^r(X, \mathcal{O}_X(-r-1)) \cong A$;
\item[(d)] the natural map
  \[
    H^0(X, \mathcal{O}_X(n)) \times H^r(X, \mathcal{O}_X(-n-r-1)) \to
    H^r(X, \mathcal{O}_X(-r-1)) \cong A
  \]
  is a perfect pairing of finitely generated free $A$-modules, for each
  $n \in \mathbf{Z}$.
\end{enumerate}
\end{theorem}

\begin{proof}
Let $\mathcal{F}$ be the quasi-coherent sheaf
$\bigoplus_{n \in \mathbf{Z}} \mathcal{O}_X(n)$. Since cohomology commutes
with arbitrary direct sums on a noetherian topological space (2.9.1), the
cohomology of $\mathcal{F}$ will be the direct sum of the cohomology of the
sheaves $\mathcal{O}(n)$. So we will compute the cohomology of $\mathcal{F}$,
and keep track
% [continued on next page]
